\section{Conditional measures}\label{sect:conditional_measures}
\subsection{Proof of Lemma \ref{lem:conditional_projection}}
First, we notice that the subspace $P\mcal{K}_{X,\pr{X}}$ consists of vectors $v\in \mcal{K}_{\mfrak{X}\backslash(X\sqcup\pr{X})}$ such that there exist $a\in \mcal{K}_{X}^{+}$, $\pr{a}\in \mcal{K}_{\pr{X}}^{-}$ and $v+a+\pr{a}\in P\mcal{K}$. Due to the property $\Gamma \mcal{K}_{X}^{\pm}=\mcal{K}_{X}^{\mp}$, we can see that
\begin{equation*}
	\Gamma_{\mfrak{X}\backslash(X\sqcup\pr{X})} P\mcal{K}_{X,\pr{X}}=\big(\overline{P}\mcal{K}+\big(\mcal{K}_{X}^{-}\oplus \mcal{K}_{\pr{X}}^{+}\big)\big)\cap \mcal{K}_{\mfrak{X}\backslash (X\cup \pr{X})}=\overline{P}\mcal{K}_{\pr{X},X}.
\end{equation*}
We can also see that $ (P\mcal{K}_{X,\pr{X}})^{\perp}=\overline{P}\mcal{K}_{\pr{X},X}$. In fact, for $u\in P\mcal{K}_{X,\pr{X}}$ and $v\in \overline{P}\mcal{K}_{\pr{X},X}$, we can take $a\in \mcal{K}_{X}^{+}$, $\pr{a}\in \mcal{K}_{\pr{X}}^{-}$, $b\in \mcal{K}_{X}^{-}$, and $\pr{b}\in \mcal{K}_{\pr{X}}^{+}$ such that
\begin{gather*}
	u+a+\pr{a} \in P\mcal{K}, \qquad v+b+\pr{b}\in \overline{P}\mcal{K}.
\end{gather*}
Now, since $a$, $\pr{a}$, $b$, $\pr{b}$ are mutually orthogonal and orthogonal to $u$ and $v$, we have
\begin{equation*}
	(u,v)_{\mcal{K}_{\mfrak{X}\backslash (X\cup\pr{X})}}=(u+a+\pr{a},v+b+\pr{b})_{\mcal{K}}=0,
\end{equation*}
which implies $(P\mcal{K}_{X,\pr{X}})^{\perp}=\overline{P}\mcal{K}_{\pr{X},X}$.
Therefore, $\overline{P}_{X,\pr{X}}=1-P_{X,\pr{X}}$ holds.

\subsection{Conditioning on the CAR algebra}\label{subsect:conditioning_CAR_alg}
For finite disjoint subsets $X, \pr{X}\subset\mfrak{X}$, the characteristic function on the corresponding cylinder set $C(X,\pr{X})$ is identified as
\begin{equation*}
	\chi_{C(X,\pr{X})}=\prod_{x\in X}a_{x}^{\ast}a_{x} \prod_{x\in \pr{X}} a_{x}a_{x}^{\ast}\in C(\Omega)\subset \mcal{C}(\mcal{K},\Gamma).
\end{equation*}
In general, given a state $\varphi$ over $\mcal{C}(\mcal{K},\Gamma)$ its conditioning over $C(X,\pr{X})$ is a state defined by (see, e.g.,~\cite{RedeiSummers2006})
\begin{equation*}
	\varphi (A|\chi_{C(X,\pr{X})})=\frac{\varphi (\chi_{C(X,\pr{X})}A\chi_{C(X,\pr{X})})}{\varphi (\chi_{C(X,\pr{X})})},\qquad A\in \mcal{C}(\mcal{K},\Gamma).
\end{equation*}
According to the decomposition $\mfrak{X}=X\sqcup \pr{X}\sqcup \mfrak{X}\backslash (X\sqcup\pr{X})$, we have a decomposition
\begin{equation*}
	\mcal{C}(\mcal{K},\Gamma)=\mcal{C}(\mcal{K}_{X\sqcup \pr{X}},\Gamma_{X\sqcup\pr{X}})\otimes \mcal{C}(\mcal{K}_{\mfrak{X}\backslash (X\sqcup\pr{X})}),
\end{equation*}
which enables us to regard $C(\mcal{K}_{\mfrak{X}\backslash (X\sqcup\pr{X})},\Gamma_{\mfrak{X}\backslash (X\sqcup\pr{X})})$ as a subalgebra of $\mcal{C}(\mcal{K},\Gamma)$ that coincides with the subalgebra realized as $\chi_{C(X,\pr{X})}\mcal{C}(\mcal{K},\Gamma)\chi_{C(X,\pr{X})}$.
Therefore, the conditional state $\varphi (\cdot|\chi_{C(X,\pr{X})})$ is regarded as a state over $C(\mcal{K}_{\mfrak{X}\backslash (X\sqcup\pr{X})},\Gamma_{\mfrak{X}\backslash (X\sqcup\pr{X})})$ and computed as
\begin{equation*}
	\varphi (A|\chi_{C(X,\pr{X})})=\frac{\varphi (A\chi_{C(X,\pr{X})})}{\varphi (\chi_{C(X,\pr{X})})},\qquad A\in C(\mcal{K}_{\mfrak{X}\backslash (X\sqcup\pr{X})},\Gamma_{\mfrak{X}\backslash (X\sqcup\pr{X})}).
\end{equation*}
In particular, in the case when $\varphi=\varphi_{S}$ is a quasi-free state associated with $S\in \mcal{Q}(\mcal{K},\Gamma)$, we have
\begin{equation*}
	\varphi_{S}(F|\chi_{C(X,\pr{X})})=\int_{\Omega (\mfrak{X}\backslash (X\sqcup\pr{X}))}F(\omega)M^{S}_{X,\pr{X}}({\rm d}\omega),\qquad F\in C(\Omega (\mfrak{X}\backslash (X\sqcup\pr{X}))).
\end{equation*}

\subsection{Some observations}
We see that obtaining $P_{X,\pr{X}}$ from a projection operator $P\in\mrm{Gr}(\mcal{K},\Gamma)$ is decomposed into fundamental steps.
\begin{Lemma}
\label{lem:rest_decomp} We have the following decomposition properties.
\begin{enumerate}\itemsep=0pt
\item[$1.$] 	When $X, \pr{X}\subset\mfrak{X}$ are disjoint finite subsets, we have
		\begin{equation}		\label{eq:rest_proj_decomp}
			P_{X,\pr{X}}=(P_{\varnothing,\pr{X}})_{X,\varnothing}=(P_{X,\varnothing})_{\varnothing, \pr{X}}.
		\end{equation}
\item[$2.$] 	When $X=X_{1}\sqcup X_{2}\subset \mfrak{X}$ is a disjoint union of finite subsets, then
		\begin{equation*}
			P_{X,\varnothing}=(P_{X_{1},\varnothing})_{X_{2},\varnothing}.
		\end{equation*}
\item[$3.$] 	When $\pr{X}=\pr{X}_{1}\sqcup \pr{X}_{2}\subset\mfrak{X}$ is a disjoint union of finite subsets, then
		\begin{equation*}
			P_{\varnothing,\pr{X}}=(P_{\varnothing,\pr{X}_{1}})_{\varnothing,\pr{X}_{2}}.
		\end{equation*}
\end{enumerate}
\end{Lemma}
\begin{proof}We only prove (\ref{eq:rest_proj_decomp}) since the other two follow from arguments of the same type.
We have noticed that $P\mcal{K}$ consists of vectors $u\in \mcal{K}_{\mfrak{X}\backslash (X\sqcup\pr{X})}$ such that there exist $a\in \mcal{K}_{X}^{+}$, $\pr{a}\in \mcal{K}_{\pr{X}}^{-}$ and $u+a+\pr{a}\in P\mcal{K}$. This exactly means that
\begin{equation*}
	u+a\in \{v\in \mcal{K}_{\mfrak{X}\backslash \pr{X}}\,|\, \exists\, \pr{a}\in \mcal{K}_{\pr{X}},\, v+\pr{a}\in P\mcal{K}\},
\end{equation*}
while the latter set is just $P\mcal{K}_{\varnothing,\pr{X}}$. Therefore, $P\mcal{K}_{X,\pr{X}}=(P\mcal{K}_{\varnothing,\pr{X}})_{X,\varnothing}$. The other relation $P\mcal{K}_{X,\pr{X}}=(P\mcal{K}_{X,\varnothing})_{\varnothing,\pr{X}}$ is also shown. Consequently, we reach~(\ref{eq:rest_proj_decomp}).
\end{proof}

We also notice that the regularity is inherited along decomposition.
\begin{Lemma}\label{lem:regularity_decomp} We have the following properties regarding the regularity.
\begin{enumerate}\itemsep=0pt
\item[$1.$] 	For disjoint finite subsets $X,\pr{X}\subset\mfrak{X}$, a projection operator $P\in\mrm{Gr}(\mcal{K},\Gamma)$ is $(X,\pr{X})$-regular if and only if one of the following equivalent conditions holds:
		\begin{enumerate}\itemsep=0pt
		\item[$(a)$] 	$P$ is $(X,\varnothing)$-regular and $P_{X,\varnothing}$ is $(\varnothing,\pr{X})$-regular.
		\item[$(b)$] 	$P$ is $(\varnothing,\pr{X})$-regular and $P_{\varnothing,\pr{X}}$ is $(X,\varnothing)$-regular.
		\end{enumerate}
\item[$2.$] 	When $X=X_{1}\sqcup X_{2}$ is a disjoint union of finite subsets, a projection operator $P\in\mrm{Gr}(\mcal{K},\Gamma)$ is $(X,\varnothing)$-regular if and only if $P$ is $(X_{1},\varnothing)$-regular and $P_{X_{1},\varnothing}$ is $(X_{2},\varnothing)$-regular.
\item[$3.$] 	When $\pr{X}=\pr{X}_{1}\sqcup \pr{X}_{2}$ is a disjoint union of finite subsets, a projection operator $P\in\mrm{Gr}(\mcal{K},\Gamma)$ is $(\varnothing,\pr{X})$-regular if and only if $P$ is $(\varnothing,\pr{X}_{1})$-regular and $P_{\varnothing,\pr{X}_{1}}$ is $(\varnothing,X_{2})$-regular.
\end{enumerate}
\end{Lemma}
\begin{proof}
We only prove the equivalence to the item (a) in (1) since the others are shown in similar arguments.
Suppose that $P\in \mrm{Gr}(\mcal{K},\Gamma)$ is $(X,\pr{X})$-regular. Then, it is obvious that $P$ is $(X,\varnothing)$-regular. Now, $P\mcal{K}_{X,\varnothing}$ consists of vectors $u\in\mcal{K}_{\mfrak{X}\backslash X}$ such that there exists $a\in \mcal{K}_{X}^{+}$ and $u+a\in P\mcal{K}$. Let us take a vector $u_{0}\in P\mcal{K}_{X,\varnothing}\cap \mcal{K}_{\pr{X}}^{-}$. Then, there exists $a_{0}\in \mcal{K}_{X}^{+}$ such that $u_{0}+a_{0}\in P\mcal{K}$, while we also have $u_{0}+a_{0}\in \mcal{K}_{X}^{+}\oplus \mcal{K}_{\pr{X}}^{-}$, which implies that $u_{0}+a_{0}\in P\mcal{K}\cap \mcal{K}_{X}^{+}\oplus \mcal{K}_{\pr{X}}^{-}=\{0\}$ by assumption. Therefore, we must have $u_{0}=0$ and see that $P\mcal{K}_{X,\varnothing}$ is $(\varnothing,\pr{X})$-regular.

Conversely, let us suppose that $P$ is $(X,\varnothing)$-regular and $P_{X,\varnothing}$ is $(\varnothing,\pr{X})$-regular. Take a~nonzero vector $u_{1}\in P\mcal{K}\cap \mcal{K}_{X}^{+}\oplus \mcal{K}_{\pr{X}}^{-}$. Since $P$ is $(X,\varnothing)$-regular, when we write $u_{1}=a_{1}+\pr{a}_{1}$, $a_{1}\in \mcal{K}_{X}^{+}$, $\pr{a}\in \mcal{K}_{\pr{X}}^{-}$, we must have $\pr{a}_{1}\neq 0$. On the other hand, by definition, $\pr{a}_{1}\in P\mcal{K}_{X,\varnothing}$, which implies that $\pr{a}_{1}\in P\mcal{K}_{X,\varnothing}\cap \mcal{K}_{\pr{X}}^{-}$ contradicting the assumption that $P_{X,\varnothing}$ is $(\varnothing,\pr{X})$-regular. Therefore, $P$ is $(X,\pr{X})$-regular.
\end{proof}

Lemmas \ref{lem:rest_decomp} and \ref{lem:regularity_decomp} verify that it suffices to prove Theorem \ref{thm:conditional_projection} in the case when $(X,\pr{X})=(\{x\},\varnothing)$ and $(\varnothing,\{x\})$ with $x\in\mfrak{X}$.
Let us introduce subspaces
\begin{gather*}
	\mcal{R}_{\pm}(x;\mcal{K}):=\big\{u\in \mcal{K}\,|\, \big(\mcal{K}^{\pm}_{\{x\}},u\big)_{\mcal{K}}=0\big\},\qquad x\in \mfrak{X},
\end{gather*}
and set $\mcal{R}_{\pm}(x;P\mcal{K})=P\mcal{K}\cap \mcal{R}_{\pm}(x;\mcal{K})$.
When we write an element $u\in P\mcal{K}$ as $u=(f,g)$ according to the direct sum decomposition $\mcal{K}=\ell^{2}(\mfrak{X})\oplus \ell^{2}(\mfrak{X})$, these are just
\begin{gather*}
	\mcal{R}_{+}(x;P\mcal{K}) =\{(f,g)\in P\mcal{K}\,|\,f(x)=0\}, \qquad \mcal{R}_{-}(x;P\mcal{K}) =\{(f,g)\in P\mcal{K}\,|\,g(x)=0\}.
\end{gather*}
Then, we can see that $P\mcal{K}_{\{x\},\varnothing}$ is the image of $\mcal{R}_{-}(x;P\mcal{K})$ under the orthogonal projection onto~$\mcal{K}_{\mfrak{X}\backslash\{x\}}$, and similarly, $P\mcal{K}_{\varnothing,\{x\}}$ is the image of $\mcal{R}_{+}(x;P\mcal{K})$ under the orthogonal projection onto~$\mcal{K}_{\mfrak{X}\backslash\{x\}}$.

\subsection[Computation of P(x,0) and P(0,x)]{Computation of $\boldsymbol{P_{\{x\},\varnothing}}$ and $\boldsymbol{P_{\varnothing,\{x\}}}$}

We introduce a lemma found in \cite[Section~7]{BufetovOlshanski2019} that plays prominent roles in computation of~$P_{\{x\},\varnothing}$ and~$P_{\varnothing,\{x\}}$.
\begin{Lemma}[\cite{BufetovOlshanski2019}]\label{lem:projection_reduction_formula}
Let $\mcal{H}$ be a Hilbert space, $\mcal{H}=\mcal{H}_{1}\oplus \mcal{H}_{2}$ be its direct sum decomposition such that $\dim \mcal{H}_{2}<\infty$ and write $\pi_{1}\colon \mcal{H}\to\mcal{H}_{1}$ for the canonical projection. Let $\mcal{L}\subset \mcal{H}$ be a~closed subspace and write the orthogonal projection $P$ onto $\mcal{L}$ in the form of
\begin{equation*}
	P=\left(
	\begin{matrix}
	A & B \\
	C & D
	\end{matrix}
	\right)
\end{equation*}
according to the direct sum decomposition $\mcal{H}=\mcal{H}_{1}\oplus \mcal{H}_{2}$, where $A\in \mfrak{B}(\mcal{H}_{1})$, $B\in \mfrak{B}(\mcal{H}_{2},\mcal{H}_{1})$, $C\in \mfrak{B}(\mcal{H}_{1},\mcal{H}_{2})$, and $D\in \mfrak{B}(\mcal{H}_{2})$.
\begin{enumerate}\itemsep=0pt
\item[$1.$] 	We define a subspace $\mcal{L}_{1}:=\mcal{L}\cap \mcal{H}_{1}$ of $\mcal{H}_{1}$ and write $P_{1}$ for the orthogonal projection onto~$\mcal{L}_{1}$ in~$\mcal{H}_{1}$. If~$D$ is invertible, i.e., $\mcal{L}^{\perp}\cap\mcal{H}_{2}=\{0\}$, we have $P_{1}=A-BD^{-1}C$.
\item[$2.$] 	We define a subspace $\widetilde{\mcal{L}}_{1}:=\pi_{1}(\mcal{L})$ of $\mcal{H}_{1}$ and write $\widetilde{P}_{1}$ for the orthogonal projection onto~$\widetilde{\mcal{L}}_{1}$ in $\mcal{H}_{1}$. If $1-D$ is invertible, i.e., $\mcal{L}\cap\mcal{H}_{2}=\{0\}$, we have $\widetilde{P}_{1}=A+B(1-D)^{-1}C$.
\end{enumerate}
\end{Lemma}
\begin{proof}
Since $P$ is an orthogonal projection, we have $A^{\ast}=A$, $B^{\ast}=C$, $D^{\ast}=D$ and
\begin{gather}\label{eq:relation_component_peojection}
	A^{2}+BC =A, \qquad AB+BD =B, \qquad CA+DC = C, \qquad CB+D^{2} =D.
\end{gather}

1.~Set $\pr{P}_{1}:=A-BD^{-1}C$, which is obviously self-adjoint and is also checked by means of~(\ref{eq:relation_component_peojection}) to be an idempotent.
		Therefore, $\pr{P}_{1}$ is the orthogonal projection onto a closed subspace $\pr{\mcal{L}}_{1}\subset\mcal{H}_{1}$.
		Suppose that $\xi\in\mcal{L}_{1}$. Then $(\xi,0)\in\mcal{H}_{1}\oplus \mcal{H}_{2}$ belongs to $\mcal{L}$ and fixed by $P$:
		\begin{equation*}
		\left(
		\begin{matrix}
		A & B \\
		C & D
		\end{matrix}
		\right)\left(
		\begin{matrix}
		\xi \\
		0
		\end{matrix}
		\right)=\left(
		\begin{matrix}
		\xi \\
		0
		\end{matrix}
		\right),
		\end{equation*}
		which is equivalent to $A\xi=\xi$ and $C\xi=0$.
		Therefore, $\xi$ is fixed by $\pr{P}_{1}$ implying that $\mcal{L}_{1}\subset \pr{\mcal{L}}_{1}$.
		Conversely, suppose that $\xi\in \pr{\mcal{L}}_{1}$. Then, we have $\big(A-BD^{-1}C\big)\xi=\xi$ implying that
		\begin{equation*}
			(\xi,A\xi)-\big(\xi,BD^{-1}C\xi\big)=(\xi,\xi).
		\end{equation*}
		Since $P$ is an orthogonal projection, we have $(\xi,A\xi)\le (\xi,\xi)$. Due to the property $B^{\ast}=C$, the operator $BD^{-1}C$ is a positive operator and $\big(\xi,BD^{-1}C\xi\big)\ge 0$. Therefore, we must have $A\xi=\xi$ and $C\xi=0$ implying $\pr{\mcal{L}}_{1}\subset \mcal{L}_{1}$.

2.~The operator $A+B(1-D)^{-1}C$ is self-adjoint and shown to be an idempotent.
		An element $\xi=(\xi_{1},\xi_{2})\in\mcal{H}_{1}\oplus\mcal{H}_{2}$ belongs to $\mcal{L}$ if and only if
		\begin{equation*}
			A\xi_{1}+B\xi_{2}=\xi_{1},\qquad C\xi_{1}+D\xi_{2}=\xi_{2}.
		\end{equation*}
		Since we have assumed that $1-D$ is invertible, the latter equation gives us $\xi_{2}=(1-D)^{-1}C\xi_{1}$, substitution of which into the former one gives $\big(A+B(1-D)^{-1}C\big)\xi_{1}=\xi_{1}$. Therefore, we obtain the desired result.
\end{proof}

Now, we can obtain explicit expressions of $P_{\{x\},\varnothing}$ and $P_{\varnothing,\{x\}}$.
Let us first write $P\in \mrm{Gr}(\mcal{K},\Gamma)$ as
\begin{equation*}
	P=\left(
	\begin{matrix}
	P_{11} & P_{12} \\
	P_{21} & P_{22}
	\end{matrix}
	\right)
\end{equation*}
according to the direct sum decomposition $\mcal{K}=\ell^{2}(\mfrak{X})\oplus \ell^{2}(\mfrak{X})$, and further write
\begin{equation*}
	P_{ij}=\left(
	\begin{matrix}
		a_{ij} & b_{ij} \\
		c_{ij} & d_{ij}
	\end{matrix}
	\right),\qquad i,j=1,2,
\end{equation*}
according to the direct sum decomposition $\ell^{2}(\mfrak{X})=\ell^{2}(\mfrak{X}\backslash\{x\})\oplus \ell^{2}(\{x\})$.

\begin{Proposition}\label{prop:expression_cond_proj}
Under the above notation, we have, if $P$ is $(\{x\},\varnothing)$-regular,
\begin{equation*}
	P_{\{x\},\varnothing}=\left(
	\begin{matrix}
		a_{11}-b_{12}d_{22}^{-1}c_{21}+b_{11}(1-d_{11})^{-1}c_{11} & a_{12}-b_{12}d_{22}^{-1}c_{22}+b_{11}(1-d_{11})^{-1}c_{12}\\
		a_{21}-b_{22}d_{22}^{-1}c_{21}+b_{21}(1-d_{11})^{-1}c_{11} & a_{22}-b_{22}d_{22}^{-1}c_{22}+b_{21}(1-d_{11})^{-1}c_{12}
	\end{matrix}
	\right)
\end{equation*}
and, if $P$ is $(\varnothing,\{x\})$-regular,
\begin{equation*}
	P_{\varnothing,\{x\}}=\left(
	\begin{matrix}
		a_{11}-b_{11}d_{11}^{-1}c_{11}+b_{12}(1-d_{22})^{-1}c_{21} & a_{12}-b_{11}d_{11}^{-1}c_{12}+b_{12}(1-d_{22})^{-1}c_{22}\\
		a_{21}-b_{21}d_{11}^{-1}c_{11}+b_{22}(1-d_{22})^{-1}c_{21} & a_{22}-b_{21}d_{11}^{-1}c_{12}+b_{22}(1-d_{22})^{-1}c_{22}
	\end{matrix}
	\right)
\end{equation*}
according to the direct sum decomposition $\mcal{K}_{\mfrak{X}\backslash\{x\}}=\mcal{K}^{+}_{\mfrak{X}\backslash\{x\}}\oplus \mcal{K}^{-}_{\mfrak{X}\backslash\{x\}}$.
\end{Proposition}
\begin{proof}
As we saw in Remark \ref{rem:check_anti-sym_matrix_func}, $P_{12}$ and $P_{21}$ are anti-symmetric; in particular, $d_{12}=d_{21}=0$. We also saw that $P_{11}$ and $P_{22}$ are self-adjoint and $JP_{11}J=1-P_{22}$, which implies $d_{11}=1-d_{22}$.

Let us derive the expression of $P_{\{x\},\varnothing}$. By assumption that $P$ is $(\{x\},\varnothing)$-regular, we have $d_{11}\neq 1$ and $d_{22}\neq 0$. Therefore, the desired expression makes sense. It is convenient to write $P$ as
\begin{equation*}
	P=\left(
	\begin{matrix}
	A & b_{1} & b_{2} \\
	c_{1}^{\mrm{T}} & d_{11} & 0 \\
	c_{2}^{\mrm{T}} & 0 & d_{22}
	\end{matrix}
	\right)
\end{equation*}
according to the direct sum decomposition $\mcal{K}=\mcal{K}_{\mfrak{X}\backslash \{x\}}\oplus \mcal{K}_{\{x\}}^{+}\oplus \mcal{K}_{\{x\}}^{-}$, where
\begin{gather*}
	A =\left(
	\begin{matrix}
		a_{11} & a_{12} \\
		a_{21} & a_{22}
	\end{matrix}
	\right), \qquad
	b_{1} =\left(
	\begin{matrix}
		b_{11} \\
		b_{21}
	\end{matrix}
	\right), \qquad
	b_{2} =\left(
	\begin{matrix}
		b_{12} \\
		b_{22}
	\end{matrix}
	\right), \\
	c_{1}^{\mrm{T}} =\left(
	\begin{matrix}
		c_{11} & c_{12}
	\end{matrix}
	\right), \qquad
	c_{2}^{\mrm{T}} =\left(
	\begin{matrix}
		c_{21} & c_{22}
	\end{matrix}
	\right).
\end{gather*}
Notice that $\mcal{R}_{-}(x;\mcal{K})=\mcal{K}_{\mfrak{X}\backslash\{x\}}\oplus \mcal{K}_{\{x\}}^{+}$.
Since $d_{22}\neq 0$, we can apply the formula (1) in Lemma~\ref{lem:projection_reduction_formula} for $\mcal{H}_{1}=\mcal{R}_{-}(x;\mcal{K})$ and $\mcal{H}_{2}=\mcal{K}_{\{x\}}^{-}$ to express the orthogonal projection onto $\mcal{R}_{-}(x;P\mcal{K})$ in $\mcal{R}_{-}(x;\mcal{K})$ as
\begin{equation*}
	\left(
	\begin{matrix}
	A & b_{1} \\
	c_{1}^{\mrm{T}} & d_{11}
	\end{matrix}
	\right)-
	\left(
	\begin{matrix}
		b_{2} \\
		0
	\end{matrix}
	\right)d_{22}^{-1}\left(
	\begin{matrix}
		c_{2}^{\mrm{T}} & 0
	\end{matrix}
	\right)=
	\left(
	\begin{matrix}
		A-b_{2}d_{22}^{-1}c_{2}^{\mrm{T}} & b_{1} \\
		c_{1}^{\mrm{T}} & d_{11}
	\end{matrix}
	\right)
\end{equation*}
according to the direct sum decomposition $\mcal{R}_{-}(x;\mcal{K})=\mcal{K}_{\mfrak{X}\backslash\{x\}}\oplus \mcal{K}_{\{x\}}^{+}$.
Recall that $P\mcal{K}_{\{x\},\varnothing}$ is the image of $\mcal{R}_{-}(x;P\mcal{K})$ under the projection onto $\mcal{K}_{\mfrak{X}\backslash\{x\}}$.
Now, since $d_{11}\neq 1$, we can apply the formula~(2) of Lemma~\ref{lem:projection_reduction_formula} for $\mcal{H}_{1}=\mcal{K}_{\mfrak{X}\backslash\{x\}}$ and $\mcal{H}_{2}=\mcal{K}_{\{x\}}^{+}$ to express $P_{\{x\},\varnothing}$ as
\begin{equation*}
	P_{\{x\},\varnothing}=A-b_{2}d_{22}^{-1}c_{2}^{\mrm{T}}+b_{1}(1-d_{11})^{-1}c_{1}^{\mrm{T}},
\end{equation*}
which is exactly the desired result for $P_{\{x\},\varnothing}$.
The expression of $P_{\varnothing,\{x\}}$ can also be derived in a~similar way to prove the desired result.
\end{proof}

\subsection{Correlation functions of conditional measures}
Let us proceed to computation of the correlation functions under conditional measures.
Let us take $P\in \mrm{Gr}(\mcal{K},\Gamma)$ that is $(\{x\},\varnothing)$-regular. In this case $M^{P}(C(\{x\},\varnothing))=\mbb{K}^{M^{P}}_{12}(x,x)=(e_{x},P_{22}e_{x})$ is non-zero. Therefore, the conditional measure $M^{P}_{\{x\},\varnothing}$ on $\Omega (\mfrak{X}\backslash\{x\})$ is well-defined. Due to the arguments in Section~\ref{subsect:conditioning_CAR_alg}, a correlation function of $M^{P}_{\{x\},\varnothing}$ is computed as
\begin{gather*}
	\rho^{M^{P}_{\{x\},\varnothing}}(x_{1},\dots, x_{n}) =\frac{\varphi_{P}(a_{x_{1}}^{\ast}\cdots a_{x_{n}}^{\ast}a_{x_{n}}\cdots a_{x_{1}}a_{x}^{\ast}a_{x})}{\varphi_{P}(a_{x}^{\ast}a_{x})}
	 =\frac{\pf \left[\mbb{K}_{P}(x_{i},x_{j})\right]_{0\le i,j\le n}}{\mbb{K}^{M^{P}}_{12}(x,x)},
\end{gather*}
where we set $x_{0}:=x$ and $x_{1},\dots, x_{n}\in \mfrak{X}\backslash \{x\}$ are distinct.
It is standard to express the above quantity in terms of the Pfaffian of a matrix of smaller size.
In fact, by means of the formula $\pf (B^{\mrm{T}}AB)=(\det B)(\pf A)$ for an anti-symmetric matrix $A$ and a matrix $B$, we can see that
\begin{equation*}
	\rho^{M^{P}_{\{x\},\varnothing}}(x_{1},\dots, x_{n})=\pf \big[\widetilde{\mbb{K}}(x_{i},x_{j})\big]_{1\le i,j\le n},
\end{equation*}
where
\begin{gather*}
\widetilde{\mbb{K}}_{11}(y,z)=\mbb{K}^{M^{P}}_{11}(y,z)-\frac{\mbb{K}^{M^{P}}_{12}(y,x)\mbb{K}^{M^{P}}_{11}(x,z)}{\mbb{K}^{M^{P}}_{12}(x,x)}+\frac{\mbb{K}^{M^{P}}_{11}(y,x)\mbb{K}^{M^{P}}_{12}(x,z)}{\mbb{K}^{M^{P}}_{12}(x,x)}, \\
\widetilde{\mbb{K}}_{12}(y,z)=\mbb{K}^{M^{P}}_{12}(y,z)-\frac{\mbb{K}^{M^{P}}_{12}(y,x)\mbb{K}^{M^{P}}_{12}(x,z)}{\mbb{K}^{M^{P}}_{12}(x,x)}+\frac{\mbb{K}^{M^{P}}_{11}(y,x)\mbb{K}^{M^{P}}_{22}(x,z)}{\mbb{K}^{M^{P}}_{12}(x,x)}, \\
\widetilde{\mbb{K}}_{22}(y,z)=\mbb{K}^{M^{P}}_{22}(y,z)-\frac{\mbb{K}^{M^{P}}_{22}(y,x)\mbb{K}^{M^{P}}_{12}(x,z)}{\mbb{K}^{M^{P}}_{12}(x,x)}+\frac{\mbb{K}^{M^{P}}_{12}(y,x)\mbb{K}^{M^{P}}_{22}(x,z)}{\mbb{K}^{M^{P}}_{12}(x,x)},
\end{gather*}
for $y,z\in\mfrak{X}\backslash\{x\}$. Note that the other component $\widetilde{\mbb{K}}_{21}(y,z)$ is determined by the anti-symmetry $\widetilde{\mbb{K}}(y,z)^{\mrm{T}}=-\widetilde{\mbb{K}}(z,y)$.

We next assume that $P\in \mrm{Gr}(\mcal{K},\Gamma)$ is $(\varnothing,\{x\})$-regular. Then, $M^{P}(C(\varnothing,\{x\}))=1-\mbb{K}^{M^{P}}_{12}(x,x)=(e_{x},P_{11}e_{y})$ is non-zero. Therefore, the conditional measure $M^{P}_{\varnothing,\{x\}}$ on $\Omega (\mfrak{X}\backslash\{x\})$ is well-defined. A correlation function is computed as
\begin{gather*}
	\rho^{M^{P}_{\varnothing,\{x\}}}(x_{1},\dots, x_{n}) =\frac{\varphi_{P}(a_{x_{1}}^{\ast}\cdots a_{x_{n}}^{\ast}a_{x_{n}}\cdots a_{x_{1}}a_{x}a^{\ast}_{x})}{\varphi_{P}(a_{x}a^{\ast}_{x})}
	 =\frac{\pf [\pr{\mbb{K}}_{P}(x_{i},x_{j}) ]_{0\le i,j\le n}}{1-\mbb{K}^{M^{P}}_{12}(x,x)},
\end{gather*}
where we set $x_{0}:=x$, $x_{1},\dots, x_{n}\in \mfrak{X}\backslash\{x\}$ are distinct and
\begin{gather*}
	\pr{\mbb{K}}_{P}(x_{i},x_{j}):=
	\begin{cases}
		\mbb{K}_{P}(x_{i},x_{j}), & 1\le i,j\le n, \\
		\left(
		\begin{matrix}
			(e_{x},P_{11}e_{x_{j}}) & (e_{x},P_{12}e_{x_{j}}) \\
			(e_{x}, P_{21}e_{x_{j}}) & (e_{x},P_{22}e_{x_{j}})
		\end{matrix}
		\right), & i=0,\quad 1\le j\le n, \vspace{1mm}\\
		\left(
		\begin{matrix}
			(e_{x_{i}},P_{22}e_{x}) & (e_{x_{i}},P_{21}e_{x}) \\
			(e_{x_{i}}, P_{12}e_{x}) & (e_{x_{i}},P_{11}e_{x})
		\end{matrix}
		\right), & 1\le i \le n,\quad j=0, \vspace{1mm}\\
		\left(
		\begin{matrix}
			0 & (e_{x},P_{11}e_{x}) \\
			(e_{x}, (P_{22}-1)e_{x}) & 0
		\end{matrix}
		\right), & i=j=0.
	\end{cases}
\end{gather*}
Again, we write the correlation function in terms of the Pfaffian of a $(2n\times 2n)$-anti-symmetric matrix. Consequently, we obtain
\begin{equation*}
	\rho^{M^{P}_{\varnothing,\{x\}}}(x_{1},\dots, x_{n})=\pf\big[\what{\mbb{K}}(x_{i},x_{j})\big]_{1\le i,j\le n},
\end{equation*}
where
\begin{gather*}
\what{\mbb{K}}_{11}(y,z) =\mbb{K}^{M^{P}}_{11}(y,z)-\frac{\mbb{K}^{M^{P}}_{11}(y,x)\mbb{K}^{M^{P}}_{12}(x,z)}{1-\mbb{K}^{M^{P}}_{12}(x,x)}+\frac{\mbb{K}^{M^{P}}_{12}(y,x)\mbb{K}^{M^{P}}_{11}(x,z)}{1-\mbb{K}^{M^{P}}_{12}(x,x)}, \\
\what{\mbb{K}}_{12}(y,z) =\mbb{K}^{M^{P}}_{12}(y,z)-\frac{\mbb{K}^{M^{P}}_{11}(y,x)\mbb{K}^{M^{P}}_{22}(x,z)}{1-\mbb{K}^{M^{P}}_{12}(x,x)}+\frac{\mbb{K}^{M^{P}}_{12}(y,x)\mbb{K}^{M^{P}}_{12}(x,z)}{1-\mbb{K}^{M^{P}}_{12}(x,x)}, \\
\what{\mbb{K}}_{22}(y,z) =\mbb{K}^{M^{P}}_{22}(y,z)-\frac{\mbb{K}^{M^{P}}_{12}(y,x)\mbb{K}^{M^{P}}_{22}(x,z)}{1-\mbb{K}^{M^{P}}_{12}(x,x)}+\frac{\mbb{K}^{M^{P}}_{22}(y,x)\mbb{K}^{M^{P}}_{12}(x,z)}{1-\mbb{K}^{M^{P}}_{12}(x,x)}
\end{gather*}
for $y,z\in\mfrak{X}\backslash\{x\}$, and the other component is determined by the anti-symmetry $\what{\mbb{K}}(y,z)^{\mrm{T}}=-\what{\mbb{K}}(z,y)$.

\subsection{Proof of Theorem \ref{thm:conditional_projection}}
It remains to show that $\widetilde{\mbb{K}}(y,z)=\mbb{K}_{P_{\{x\},\varnothing}}(y,z)$ and $\what{\mbb{K}}(y,z)=\mbb{K}_{P_{\varnothing,\{x\}}}(y,z)$, $y,z\in \mfrak{X}\backslash\{x\}$.
It is straightforward that
\begin{gather*}
\widetilde{\mbb{K}}_{11}(y,z)=(e_{y},P_{21}e_{z})-\frac{(e_{y},P_{22}e_{x})(e_{x},P_{21}e_{z})}{(e_{x},P_{22}e_{x})}+\frac{(e_{y},P_{21}e_{x})(e_{x},P_{11}e_{z})}{(e_{x},(1-P_{11})e_{x})}, \\
\widetilde{\mbb{K}}_{12}(y,z)=(e_{y},P_{22}e_{z})-\frac{(e_{y},P_{22}e_{x})(e_{x},P_{22}e_{z})}{(e_{x},P_{22}e_{x})}+\frac{(e_{y},P_{21}e_{x})(e_{x},P_{12}e_{z})}{(e_{x},(1-P_{11})e_{x})}, \\
\widetilde{\mbb{K}}_{22}(y,z)=(e_{y},P_{12}e_{z})-\frac{(e_{y},P_{12}e_{x})(e_{x},P_{22}e_{z})}{(e_{x},P_{22}e_{x})}+\frac{(e_{y},P_{11}e_{x})(e_{x},P_{12}e_{z})}{(e_{x},(1-P_{11})e_{x})}.
\end{gather*}
Comparing these descriptions with the result of Proposition~\ref{prop:expression_cond_proj}, we conclude that $\widetilde{\mbb{K}}(y,z)=\mbb{K}_{P_{\{x\},\varnothing}}(y,z)$, $y,z\in\mfrak{X}\backslash\{x\}$.
It also follows from
\begin{gather*}
\what{\mbb{K}}_{11}(y,z)=(e_{y},P_{21}e_{z})-\frac{(e_{y}P_{21}e_{x})(e_{x},P_{11}e_{z})}{(e_{x},P_{11}e_{x})}+\frac{(e_{y},P_{22}e_{x})(e_{x},P_{21}e_{z})}{(e_{x},(1-P_{22})e_{x})}, \\
\what{\mbb{K}}_{12}(y,z)=(e_{y},P_{22}e_{z})-\frac{(e_{y}P_{21}e_{x})(e_{x},P_{12}e_{z})}{(e_{x},P_{11}e_{x})}+\frac{(e_{y},P_{22}e_{x})(e_{x},P_{22}e_{z})}{(e_{x},(1-P_{22})e_{x})}, \\
\what{\mbb{K}}_{22}(y,z)=(e_{y},P_{12}e_{z})-\frac{(e_{y}P_{11}e_{x})(e_{x},P_{12}e_{z})}{(e_{x},P_{11}e_{x})}+\frac{(e_{y},P_{12}e_{x})(e_{x},P_{22}e_{z})}{(e_{x},(1-P_{22})e_{x})}
\end{gather*}
and Proposition \ref{prop:expression_cond_proj} that the other desired coincidence $\what{\mbb{K}}(y,z)=\mbb{K}_{P_{\varnothing,\{x\}}}(y,z)$, $y,z\in \mfrak{X}\backslash\{x\}$ holds.

